\begin{annexitem}{annex:p08.rat2.6.hard-update}{thm:p08-rat26-hard-update}{Public restriction, knowledge and the dialogue model}
\begin{proof}
At a world where \(\phi\) is false, the left side of the reduction
law is vacuously true and the implication on the right is true. Suppose
instead that \(M,w\models\phi\). The left side says that every
\(v\sim_i w\) surviving the restriction satisfies \(\psi\) in
\(M\upharpoonright\phi\). On the right, each original accessible
\(v\) satisfies \([!\phi]\psi\): this is vacuous at removed
worlds, and requires exactly that same updated truth of \(\psi\)
at surviving worlds. Thus the conditions coincide, including arbitrary
epistemic or dynamic \(\psi\).

Restriction preserves the atomic valuations of every surviving world.
An induction on propositional construction therefore shows that factual
\(\phi\) remains true at all worlds of the restricted model. Every
world reachable from its actual world lies in that domain, so \(C\phi\)
holds. At a removed world the announcement box is again vacuous.

In the two-world example, agent 1 accesses both truth values of \(p\),
so neither \(K_1p\) nor \(K_1\neg p\) holds. Agent 2 accesses
only its own world, and hence \(K_2p\vee K_2\neg p\) holds at
both worlds. Agent 1 therefore knows that disjunction. Announcing \(p\)
at \(w_+\) restricts the model to \(\{w_+\}\); every accessibility
step remains there, proving \(Cp\). No property of a question token
was used. For the failure of the factual generalization, announce
\(p\wedge\neg K_1p\) in the initial model at \(w_+\).
The restriction is the same singleton, where \(K_1p\) now holds,
so the announced epistemic formula no longer holds.
\end{proof}
\end{annexitem}

\begin{annexitem}{annex:p08.rat2.6.soft-update}{thm:p08-rat26-soft-update}{Best-world promotion and conditional separation}
\begin{proof}
The antecedent \(\neg K\neg\phi\) is precisely nonemptiness of
\(P=\llbracket\phi\rrbracket\). Finiteness and totality of the
preorder ensure that \(\min P\) is nonempty. Conservative upgrade
makes that set the best tier of the whole model. Radical upgrade puts
all of \(P\) before its complement and preserves its internal order,
so its best tier is also \(\min P\). All worlds of that tier satisfy
\(\phi\). Because \(\phi\) is factual and atomic valuations do
not change, both updated models satisfy \(B\phi\). The domain and
factual truth values are unchanged, so knowledge of any factual
proposition is unchanged. No claim is made for epistemic formulas
whose truth depends on the ordering or on belief.

In the example, \(\min P=\{b\}\). Promotion of that set puts
\(b\) first while retaining \(a\prec c\). Radical promotion of
all of \(P\) retains \(b\prec c\) and places \(a\) last.
The best \(\chi\)-world is therefore \(a\) in the former model
and \(c\) in the latter. The valuation of \(q\) gives the stated
opposite conditional beliefs. In both models the best world is \(b\),
which satisfies \(\phi\), while world \(a\) remains possible and
falsifies it. Hence both have \(B\phi\) and \(\neg K\phi\).
\end{proof}
\end{annexitem}

\begin{annexitem}{annex:p08.rat2.6.sabotage-space}{thm:p08-rat26-sabotage-space}{Depth-first strategy evaluation with explicit edges}
\begin{proof}
Let \(m\) be the number of initially listed edge occurrences and
\(n\) the number of vertices. A position is recorded by its current
vertex, whose turn it is, and an \(m\)-bit vector of surviving edges.
The target is tested before requiring a move. At a Runner position,
the position is winning exactly when it is a target or some legal
successor is winning. At a nontarget Blocker position, it is winning
for Runner exactly when every legal deletion successor is winning.
A nontarget Runner position with no successor is losing.

Each completed round deletes one edge occurrence. Thus a play has at
most \(m\) deletion moves and at most \(m+1\) Runner moves before
termination. In fact a legal Blocker turn always has an available edge,
because Runner has just traversed one. The recursion has depth at most
\(2m+1\), and strictly exhausts edges between successive rounds;
it cannot run forever. Backward induction over this finite play tree
justifies the existential and universal tests: Runner chooses its move,
whereas a winning strategy must survive every Blocker choice.

Evaluate the tree depth first, retaining only the stack of positions,
the successor index currently explored and the partial Boolean result.
Each stack entry uses \(O(m+\log n+\log(m+n))\) bits. There are
\(O(m+1)\) entries, so total space is polynomial in the explicit
input length. Time and the size of a written strategy tree can be
exponential. The argument would not justify the same bound in the
bit-length of a succinct multiplicity encoding: its depth bound uses
the actual listed edge count. It proves membership, not hardness.
\end{proof}
\end{annexitem}

\begin{annexitem}{annex:p08.rat2.6.hawk-dove}{thm:p08-rat26-hawk-dove}{The invasion inequality and all parameter cases}
\begin{proof}
Against Hawk frequency \(z\), pure Hawk earns
\(z(V-C)/2+(1-z)V\), while pure Dove earns \((1-z)V/2\).
Their difference is \((V-Cz)/2\). Linearity in the row mixture gives
\[
 u(x,z)-u(y,z)=(x-y)(V-Cz)/2.
\]
If \(0<V<C\), put \(x^*=V/C\in(0,1)\). Against the
contaminated population \(z=(1-\varepsilon)x^*+\varepsilon y\),
\(V-Cz=C\varepsilon(x^*-y)\). Thus
\[
 u(x^*,z)-u(y,z)=C\varepsilon(x^*-y)^2/2>0
\]
for every distinct mutant and \(0<\varepsilon<1\). This proves
the invasion criterion, not merely equality of payoffs at \(x^*\).
For a resident \(x<x^*\), mutant \(y=1\) has strictly higher
payoff when \(\varepsilon\) is sufficiently small: at zero
contamination \((x-1)(V-Cx)<0\), and continuity preserves that
sign. For \(x>x^*\), mutant \(y=0\) similarly defeats the
resident, since \(x(V-Cx)<0\) initially. Hence no other resident
is stable.

If \(V\ge C\), take resident \(x=1\) and any mutant \(y<1\).
Its advantage is
\[
 (1-y)\bigl(V-C+C\varepsilon(1-y)\bigr)/2>0
\]
for every \(0<\varepsilon<1\). This remains strictly positive
at \(V=C\) because the second-order invasion term is positive.
If a resident has \(x<1\), then \(V-Cx>0\), including when
\(V=C\); the mutant \(y=1\) has a strictly positive advantage
for sufficiently small contamination. This proves uniqueness also
in the boundary case. All inequalities concern the specified payoff
model and independent pairing, not an empirical evolutionary trajectory.
\end{proof}
\end{annexitem}

\begin{annexitem}{annex:p08.rat2.6.network}{thm:p08-rat26-network}{The four transitions, eventual periodicity and rigidity}
\begin{proof}
For any active set \(X\), the indicator of the updated set is
\[
 (\mathbf1_{2\in X}\mathbf1_{3\in X},
   \mathbf1_{4\in X},\mathbf1_{2\in X},\mathbf1_{3\in X}).
\]
Substitution gives \(T(\{1\})=\varnothing\) and
\(T(\varnothing)=\varnothing\). For singleton cycle vertices,
\(T(\{2\})=\{3\}\), \(T(\{3\})=\{4\}\), and
\(T(\{4\})=\{2\}\). The set \(\{1,2\}\) also maps to
\(\{3\}\), proving its transient and the same following cycle.
Finally,
\[
 \{1,2,3\}\mapsto\{1,3,4\}\mapsto\{2,4\}
 \mapsto\{2,3\}\mapsto\{1,3,4\}.
\]
Each equality follows from the four displayed membership conditions.

On an \(n\)-vertex graph there are only \(2^n\) active sets.
Among the first \(2^n+1\) states of a deterministic trajectory,
two are equal. Applying the same function to equal states gives equal
successors; induction gives repetition with the same positive period
forever. The singleton three-cycle disproves necessary stabilization.

An automorphism of the full directed graph fixes vertex 1, the unique
vertex of outdegree two. Its outneighbor set \(\{2,3\}\) is then
preserved, and vertex 4 is the only remaining vertex distinct from 1
outside that set, so 4 is fixed. Its unique outneighbor 3 is fixed;
the unique outneighbor 2 of 3 is then fixed. Thus the automorphism
is the identity. An isomorphism between two marked copies of this
same graph would be an automorphism taking one marked set to the
other, which is impossible for distinct singletons. In contrast the
subgraph induced by each active singleton has one vertex and no
self-loop; these reduced structures are isomorphic. The comparison
must specify whether the inactive part of the graph is retained.
\end{proof}
\end{annexitem}

\begin{annexitem}{annex:p08.rat2.6.learning-limits}{thm:p08-rat26-learning-limits}{Optimization and held-out countercases}
\begin{proof}
Since \(J'(t)=2t\), gradient descent with fixed positive step
\(\eta\) gives \(t_{k+1}=(1-2\eta)t_k\), hence
\(t_k=(1-2\eta)^k t_0\). For \(\eta>1\) the multiplier
has absolute value greater than one, so nonzero initial values
diverge in magnitude. For \(\eta=1\) they alternate between
\(t_0\) and \(-t_0\), and do not converge. For comparison,
\(0<\eta<1\) gives convergence to zero. The result therefore
identifies a missing step-size condition, not failure of all gradient
procedures.

For every real \(c\), \(f_c(0)=0\), so its squared error on
the single training example is zero. But \(f_c(1)=c\), so its
squared error on the separately specified held-out example is \(c^2\).
As \(c\) varies this is unbounded. Thus the training observation
does not identify the held-out prediction without further restrictions
or inductive assumptions. This finite mathematical countercase is
not a measured result about a neural-network benchmark.
\end{proof}
\end{annexitem}
